\section*{Capítulo \ref{ch:g6:cell-unit-of-life} --- A célula, unidade da vida}

\begin{solution}{exo:g6:cell-unit-of-life:1}
A \omterm{def:g6:cell-unit-of-life:parts}{membrana} --- o envoltório fino que mantém a \omterm{def:g5:first-look-at-cells:cell}{célula} unida e controla
o que entra e o que sai; o \omterm{def:g6:cell-unit-of-life:parts}{citoplasma} --- o recheio gelatinoso onde o
trabalho acontece; o \omterm{def:g6:cell-unit-of-life:parts}{núcleo} --- o centro de comando que dirige a vida
e a divisão da \omterm{def:g5:first-look-at-cells:cell}{célula}.
\end{solution}

\begin{solution}{exo:g6:cell-unit-of-life:2}
A \omterm{prop:g6:cell-unit-of-life:plantextras}{parede celular} rígida --- firmeza de \omterm{def:g1:plants-around-us:plant}{planta} sem \omterm{def:g3:bones-and-muscles:skeleton}{ossos}; os grãos
verdes --- as oficinas que captam \omterm{def:g6:exploring-our-environment:conditions}{luz} do ofício das produtoras; a bolsa
grande de água --- a pressão dela mantém a \omterm{def:g5:first-look-at-cells:cell}{célula} túrgida.
\end{solution}

\begin{solution}{exo:g6:cell-unit-of-life:3}
Cebola: paredes retas que se encontram em quinas, com o \omterm{def:g6:cell-unit-of-life:parts}{núcleo}
empurrado para o lado. Bochecha: contorno macio e arredondado só com a
\omterm{def:g6:cell-unit-of-life:parts}{membrana}, com o \omterm{def:g6:cell-unit-of-life:parts}{núcleo} perto do meio.
\end{solution}

\begin{solution}{exo:g6:cell-unit-of-life:4}
Um organismo completo de exatamente uma \omterm{def:g5:first-look-at-cells:cell}{célula}. O \omterm{ex:g6:producing-our-food:bread}{fermento}; a remadora
em forma de chinelo da gota de lagoa (ou os nadadores verdes dela, a
bolha que escorre, os \omterm{def:g6:producing-our-food:fermentation}{micróbios} do iogurte).
\end{solution}

\begin{solution}{exo:g6:cell-unit-of-life:5}
Todo \omterm{def:g6:exploring-our-environment:components}{ser vivo} é feito de uma ou de muitas \omterm{def:g5:first-look-at-cells:cell}{células}, as menores \omterm{prop:g5:first-look-at-cells:principle}{unidades
da vida}; e toda \omterm{def:g5:first-look-at-cells:cell}{célula} vem da divisão de uma \omterm{def:g5:first-look-at-cells:cell}{célula} que já existia.
\end{solution}

\begin{solution}{exo:g6:cell-unit-of-life:6}
Dividindo-se em dois: uma \omterm{def:g5:first-look-at-cells:cell}{célula} vira dois organismos completos.
\end{solution}

\begin{solution}{exo:g6:cell-unit-of-life:7}
As \omterm{prop:g6:cell-unit-of-life:plantextras}{paredes celulares} rígidas e as bolsas de água túrgidas: milhões de
\omterm{def:g5:first-look-at-cells:cell}{células} com parede e sob pressão se escoram umas nas outras como uma
alvenaria inflada. Murchar é as bolsas ficando moles --- a pressão
perdida por falta de água, com a estrutura inteira cedendo.
\end{solution}

\begin{solution}{exo:g6:cell-unit-of-life:8}
Os grãos verdes. Os \omterm{def:g1:animals-around-us:animal}{animais} são \omterm{def:g5:food-webs:roles}{consumidores} --- transformadores da
\omterm{def:g6:plants-are-producers:matter}{matéria orgânica} que comem, e nunca fabricantes a partir do mineral
--- então as \omterm{def:g5:first-look-at-cells:cell}{células} deles não carregam grãos de fabricação: o
equipamento combina com o ofício, \omterm{def:g5:first-look-at-cells:cell}{célula} por \omterm{def:g5:first-look-at-cells:cell}{célula}.
\end{solution}

\begin{solution}{exo:g6:cell-unit-of-life:9}
Por exemplo: \omterm{def:g5:first-look-at-cells:cell}{células} de revestimento (as da bochecha), que cobrem
superfícies; \omterm{def:g5:first-look-at-cells:cell}{células} musculares, que se contraem para puxar \omterm{def:g3:bones-and-muscles:skeleton}{ossos};
\omterm{def:g5:first-look-at-cells:cell}{células} nervosas, que levam os informes e as ordens do circuito do
\omterm{def:g5:brain-in-command:system}{cérebro}.
\end{solution}

\begin{solution}{exo:g6:cell-unit-of-life:10}
Arranhão cicatrizando: \omterm{def:g5:first-look-at-cells:cell}{células} da \omterm{def:g1:the-five-senses:sense}{pele} se dividindo para reconstruir o
pedaço. Massa dobrando: \omterm{def:g5:first-look-at-cells:cell}{células} de \omterm{ex:g6:producing-our-food:bread}{fermento} se dividindo de novo e de
novo na massa quentinha e cheia de açúcar. Girino crescendo: as
\omterm{def:g5:first-look-at-cells:cell}{células} dele se dividindo e se especializando enquanto o corpo se
constrói.
\end{solution}

\begin{solution}{exo:g6:cell-unit-of-life:11}
O ofício das produtoras: fabricação, numa \omterm{def:g5:first-look-at-cells:cell}{célula} só, de \omterm{def:g6:plants-are-producers:matter}{matéria
orgânica} a partir de suprimentos minerais e de \omterm{def:g6:exploring-our-environment:conditions}{luz}. Na teia da lagoa
eles são primeiros elos --- o andar verde que os \omterm{ex:g2:animals-grow-up:frog}{girinos} e as pulgas
d'água pastam.
\end{solution}

\begin{solution}{exo:g6:cell-unit-of-life:12}
Cada uma das suas \omterm{def:g5:first-look-at-cells:cell}{células} veio de uma divisão, cuja célula-mãe veio de
uma divisão, sem uma única interrupção --- passando pelo óvulo
fecundado, pelas \omterm{def:g5:first-look-at-cells:cell}{células} dos seus pais e mais para trás, até o passado
profundo da vida. A consequência numa frase: toda \omterm{def:g5:first-look-at-cells:cell}{célula} viva, a sua
inclusive, é o elo mais recente de uma corrente ininterrupta de
divisões tão antiga quanto a própria vida.
\end{solution}

\begin{solution}{pb:g6:cell-unit-of-life:1}
\textbf{1.} A parede: a \omterm{prop:g6:cell-unit-of-life:plantextras}{parede celular} vegetal. O pontinho: o \omterm{def:g6:cell-unit-of-life:parts}{núcleo}.
O culpado pelo aperto: a bolsa grande de água, empurrando o \omterm{def:g6:cell-unit-of-life:parts}{núcleo}
para o lado.

\textbf{2.} A tem a parede, a bolsa de água (e, nas partes verdes, os
grãos verdes) que B não tem. As duas compartilham o kit: \omterm{def:g6:cell-unit-of-life:parts}{membrana},
\omterm{def:g6:cell-unit-of-life:parts}{citoplasma}, \omterm{def:g6:cell-unit-of-life:parts}{núcleo}.

\textbf{3.} Os grãos são as oficinas que captam \omterm{def:g6:exploring-our-environment:conditions}{luz} --- e um \omterm{ex:g4:plants-without-seeds:bulbtuber}{bulbo}
subterrâneo vive no escuro, onde as oficinas ficariam paradas: a
\omterm{def:g1:plants-around-us:plant}{planta} constrói esses grãos nas \omterm{def:g1:plants-around-us:parts}{folhas} iluminadas dela, e não na
despensa enterrada.

\textbf{4.} As \omterm{def:g5:first-look-at-cells:cell}{células} dos próprios alunos. Pela frase 2, cada uma
veio da divisão de uma \omterm{def:g5:first-look-at-cells:cell}{célula} que já existia --- uma corrente de
divisões que remonta ao óvulo fecundado de cada aluno.

\textbf{5.} A remadora: se move, se dirige, se alimenta --- um
organismo \omterm{def:g6:cell-unit-of-life:unicellular}{unicelular} completo (do tipo \omterm{def:g1:animals-around-us:animal}{animal}). Os nadadores verdes:
se movem e captam \omterm{def:g6:exploring-our-environment:conditions}{luz} --- produtoras \omterm{def:g6:cell-unit-of-life:unicellular}{unicelulares}. A bolha: escorre e
engole o alimento --- um caçador \omterm{def:g6:cell-unit-of-life:unicellular}{unicelular}. Cada um mostra as
atividades de uma vida inteira numa \omterm{def:g5:first-look-at-cells:cell}{célula} só.

\textbf{6.} Os grãos verdes --- o equipamento de captar \omterm{def:g6:exploring-our-environment:conditions}{luz} que toca a
fabricação das produtoras.

\textbf{7.} Elas se alimentaram do açúcar do \omterm{def:g2:animals-grow-up:born}{leite} e liberaram o ácido
que o coalhou em iogurte: são \omterm{def:g6:exploring-our-environment:components}{seres vivos} \omterm{def:g6:cell-unit-of-life:unicellular}{unicelulares} --- a força de
trabalho de \omterm{def:g6:producing-our-food:fermentation}{micróbios}, vista em pessoa enfim.

\textbf{8.} Observar um pontinho se dividir: um virando dois, dois
virando quatro --- a divisão é a assinatura que a frase 2 exige dos
vivos. (Numa lâmina aquecida, os \omterm{def:g6:producing-our-food:fermentation}{micróbios} do iogurte fazem a gentileza
em poucas horas.)

\textbf{9.} A: muitas \omterm{def:g5:first-look-at-cells:cell}{células}, vegetal. B: muitas \omterm{def:g5:first-look-at-cells:cell}{células} (uma amostra
de um \omterm{def:g1:animals-around-us:animal}{animal} multicelular), \omterm{def:g1:animals-around-us:animal}{animal}. C: \omterm{def:g5:first-look-at-cells:cell}{células} sozinhas --- do tipo
\omterm{def:g1:animals-around-us:animal}{animal} (a remadora, a bolha) e do tipo vegetal (os nadadores verdes)
lado a lado. D: \omterm{def:g5:first-look-at-cells:cell}{células} sozinhas, nem vegetais nem \omterm{def:g1:animals-around-us:animal}{animais} ---
\omterm{def:g6:producing-our-food:fermentation}{micróbios}.

\textbf{10.} \omterm{def:g6:cell-unit-of-life:parts}{Membrana}, \omterm{def:g6:cell-unit-of-life:parts}{citoplasma}, \omterm{def:g6:cell-unit-of-life:parts}{núcleo}. A universalidade dele na
cebola, na bochecha, na lagoa e no iogurte é prova da \omterm{prop:g5:first-look-at-cells:principle}{unidade da vida}:
um único plano de construção embaixo de tudo isso.

\textbf{11.} A cebola cresceu de uma \omterm{def:g5:first-look-at-cells:cell}{célula} fabricada pela planta-mãe
dela, por divisões; o aluno começou como um óvulo fecundado; cada
corrente do iogurte começou com um \omterm{def:g1:healthy-habits:germs}{micróbio} se dividindo e se
dividindo. Tudo o que foi desenhado remonta a uma única \omterm{def:g5:first-look-at-cells:cell}{célula}
inicial.

\textbf{12.} Frase 1 --- todo \omterm{def:g6:exploring-our-environment:components}{ser vivo} é feito de \omterm{def:g5:first-look-at-cells:cell}{células}: mostrada
melhor por A e B (e pela visita como um todo). Frase 2 --- toda \omterm{def:g5:first-look-at-cells:cell}{célula}
vem de uma \omterm{def:g5:first-look-at-cells:cell}{célula} que se divide: mostrada melhor por C e D, em que a
divisão é um modo de vida visível.
\end{solution}
